% Source PDF: source/普通高考/2016/2016全国1理(河南,河北,山西,江西,湖北,湖南,广东,安徽,福建).pdf
% Page: 2; PDF embedded bitmap attempt: no image objects were present (pdfimages -list).
% Grid evidence: tmp/redraw_2016_national1/page-2_grid100.png and q19_block_grid50.png.
% Original-side crop: tmp/redraw_2016_national1/q19_orig_crop.png, cropped from the no-grid page render.
% Elements: frequency y-axis, broken x-axis, four bars centered at 8, 9, 10, 11, and dashed guides at 20 and 40.
% Relations: the source graph gives frequencies 20, 40, 20, 20 for replacement counts 8, 9, 10, 11.
% solid segments: bar outlines, x/y axes, tick marks, and x-axis break stroke.
% dashed segments: horizontal reference guides at frequencies 20 and 40.
% Labels: 频数 is above the y-axis; 更换的易损零件数 is right of the x-axis arrow; tick labels stay below the axis.
\begin{tikzpicture}[baseline]
\begin{axis}[
    width=7.4cm,
    height=4.8cm,
    axis lines=none,
    clip=false,
    % Internal display coordinates compress the broken 0-to-8 segment to
    % exactly one class width; the four unit-width bars occupy [1,5).
    xmin=-0.8,
    xmax=6.8,
    ymin=-5.5,
    ymax=53,
    ticks=none,
    x tick label as interval=false,
]
    \draw[dashed,line width=0.55pt] (axis cs:0,20) -- (axis cs:5,20);
    \draw[dashed,line width=0.55pt] (axis cs:0,40) -- (axis cs:3,40);

    \addplot[
        ybar interval,
        draw=black,
        fill=white,
        line width=0.7pt,
    ] coordinates {
        (1,20)
        (2,40)
        (3,20)
        (4,20)
        (5,0)
    };

    \draw[->,line width=0.8pt] (axis cs:0,0) -- (axis cs:6.0,0);
    \draw[->,line width=0.8pt] (axis cs:0,0) -- (axis cs:0,50.5);
    \fill[white] (axis cs:0.28,-2.0) -- (axis cs:0.46,2.4) -- (axis cs:0.64,-2.0)
        -- (axis cs:0.82,2.4) -- (axis cs:1.00,-2.0) -- cycle;
    \draw[line width=0.8pt] (axis cs:0.28,-2.0) -- (axis cs:0.46,2.4)
        -- (axis cs:0.64,-2.0) -- (axis cs:0.82,2.4) -- (axis cs:1.00,-2.0);

    \node[anchor=south west,inner sep=0pt,font=\scriptsize] at (axis cs:0.25,48.0) {频数};
    \node[anchor=west,inner sep=0pt,font=\scriptsize] at (axis cs:5.15,3.0) {更换的易损零件数};
    \node[anchor=north east,inner sep=0pt,font=\scriptsize] at (axis cs:0,-1.0) {0};

    % The source chart labels the discrete values at bar centers.
    \draw[line width=0.5pt] (axis cs:1.5,0) -- (axis cs:1.5,-1.5);
    \draw[line width=0.5pt] (axis cs:2.5,0) -- (axis cs:2.5,-1.5);
    \draw[line width=0.5pt] (axis cs:3.5,0) -- (axis cs:3.5,-1.5);
    \draw[line width=0.5pt] (axis cs:4.5,0) -- (axis cs:4.5,-1.5);
    \node[anchor=north,inner sep=0pt,font=\scriptsize] at (axis cs:1.5,-2.0) {8};
    \node[anchor=north,inner sep=0pt,font=\scriptsize] at (axis cs:2.5,-2.0) {9};
    \node[anchor=north,inner sep=0pt,font=\scriptsize] at (axis cs:3.5,-2.0) {10};
    \node[anchor=north,inner sep=0pt,font=\scriptsize] at (axis cs:4.5,-2.0) {11};

    \draw[line width=0.5pt] (axis cs:0,20) -- (axis cs:-0.35,20);
    \draw[line width=0.5pt] (axis cs:0,40) -- (axis cs:-0.35,40);
    \node[anchor=east,inner sep=0pt,font=\scriptsize] at (axis cs:-0.45,20) {20};
    \node[anchor=east,inner sep=0pt,font=\scriptsize] at (axis cs:-0.45,40) {40};
\end{axis}
\end{tikzpicture}
